% !TEX root = ../main.tex
\section{Rheology, Thermal, \& Materials Physics}
The viscosity of molten glass is highly temperature-dependent, which dictates its extrudability and structural stability. This behavior is typically modeled using the Vogel-Fulcher-Tammann (VFT) equation:
\begin{equation}
\log_{10}\eta = A + \frac{B}{T - T_0}
\end{equation}
where $\eta$ is the viscosity, $T$ is the absolute temperature, and $A, B, T_0$ are empirical constants dependent on the glass composition (e.g., fused silica, borosilicate, soda-lime). 

In addition to viscosity, the gravitational slumping and surface tension must be controlled. The capillary Rayleigh-Plateau instability can cause liquid glass cords to break into droplets, meaning the extrusion rate and cooling must be precisely tuned to solidify the glass before slumping or bead breakup occurs \cite{klein2015additive}.
